% =====================================================================
%  CHƯƠNG 8. KẾT QUẢ THỰC NGHIỆM VÀ PHÂN TÍCH ĐỘNG HỌC HỌC TẬP
% =====================================================================
\chapter{Kết quả thực nghiệm và phân tích động học học tập}

\section{Ảnh chụp màn hình kết quả thực thi (Copy Screen Output)}
Các hình dưới đây là ảnh chụp màn hình thực tế kết quả chạy chương trình trên Jupyter Notebook, đúng yêu cầu của bài
tập (copy screen output): nhật ký huấn luyện trên console, lộ trình học được, hai biểu đồ trực quan và Q-Table.

\subsection{Nhật ký huấn luyện}
\anhNB{out_train}{Ảnh chụp màn hình kết quả ô 4: cấu hình bài toán và nhật ký huấn luyện sau mỗi 250 episode.}{nb-train}

\subsection{Lộ trình tối ưu}
\anhNB{out_path}{Ảnh chụp màn hình kết quả ô 5: lộ trình học được, đối chiếu với đường ngắn nhất tìm bằng
BFS.}{nb-path}

Tác tử tìm được lộ trình \SL{ch8_qlearning/result/path_len} bước với tổng phần thưởng
\SL{ch8_qlearning/result/greedy_reward}, đúng bằng độ dài đường ngắn nhất \SL{ch8_qlearning/result/bfs_len} bước do BFS
tìm trên bản đồ. Tổng phần thưởng khớp với tính toán: \SL{ch8_qlearning/result/path_len} bước gồm 23 bước thường
($-1$ mỗi bước) và bước cuối vào đích ($+100$), tức $-23 + 100 = 77$.

\subsection{Biểu đồ trực quan}
\anhNB{out_plot}{Ảnh chụp màn hình kết quả ô 7: lộ trình tối ưu trong mê cung $12 \times 12$ và đường cong hội tụ phần
thưởng sau \SL{ch8_qlearning/params/episodes} episode.}{nb-plot}

Để đọc rõ hơn, hai biểu đồ trong Hình~\ref{fig:nb-plot} được vẽ lại dạng vector ở cỡ toàn trang từ chính số liệu notebook
đã lưu (Hình~\ref{fig:ch8-route} và Hình~\ref{fig:ch8-rewards}).

\hinhpdf{ch8_route}{0.9}{Lộ trình tối ưu \SL{ch8_qlearning/result/path_len} bước của tác tử trong mê cung (vẽ lại
dạng vector).}{ch8-route}

\hinhpdf{ch8_rewards}{1}{Phần thưởng mỗi episode, trung bình trượt 100 episode và lịch giảm $\varepsilon$ (vẽ lại dạng
vector).}{ch8-rewards}

\subsection{Q-Table và lưu số liệu}
\anhNB[0.5]{out_qtable}{Ảnh chụp màn hình kết quả ô 8: trích đoạn Q-Table tại các ô then chốt của lộ trình.}{nb-qtable}
\anhNB{out_save}{Ảnh chụp màn hình ô 9: lưu số liệu thực nghiệm ra tệp JSON cho báo cáo.}{nb-save}

\section{Phân tích động học học tập và hội tụ chính sách}
Bảng~\ref{tab:ch8-log} tổng hợp nhật ký huấn luyện theo từng khối 250 episode. Kết hợp với Hình~\ref{fig:ch8-rewards},
quá trình tự học của tác tử chia thành hai pha rõ rệt.

\bang{ch8_log}{Diễn biến huấn luyện theo từng khối 250 episode.}{ch8-log}

\subsection{Pha thăm dò (khoảng 500 episode đầu)}
Với $\varepsilon$ gần 1, tác tử gần như đi ngẫu nhiên. Nó liên tục đâm vào tường lửa, mỗi lần mất 10 điểm và đứng yên,
nên phần thưởng rất thấp: trung bình 100 episode đầu chỉ \SL{ch8_qlearning/result/reward_first100}, episode tệ nhất
\SL{ch8_qlearning/result/reward_min}. Nhiều episode chạm giới hạn \SL{ch8_qlearning/params/max_steps} bước mà không về
đích; tỉ lệ về đích của 250 episode đầu là \SL{ch8_qlearning/result/success_first250Pct}\%. Lần đầu tiên tác tử tình cờ về
được đích là ở episode \SL{ch8_qlearning/result/first_success}. Từ đó phần thưởng $+100$ bắt đầu được lan ngược dần từ ô
đích về các ô xa hơn qua từng lần cập nhật Bellman. Đây là giai đoạn \emph{Problem Generator} hoạt động mạnh nhất: tác tử
tích lũy kinh nghiệm chứ chưa khai thác.

\subsection{Pha khai thác và hội tụ}
Khi giá trị lan đủ xa, các hành động tốt nổi lên trong Q-Table và $\varepsilon$ giảm làm tác tử ngày càng đi theo chúng.
Đường trung bình trượt 100 episode chuyển sang dương ở khoảng episode \SL{ch8_qlearning/result/ma_first_positive}; từ
khối 500 trở đi tác tử về đích 100\% số episode. $\varepsilon$ xuống dưới 0,05 ở episode
\SL{ch8_qlearning/result/eps_below_005}; sau đó phần thưởng ổn định gần trần 77. Trung bình 100 episode cuối đạt
\SL{ch8_qlearning/result/reward_last100} với \SL{ch8_qlearning/result/steps_last100} bước. Con số này thấp hơn 77 một
chút vì $\varepsilon$ vẫn giữ ở mức tối thiểu \SL{ch8_qlearning/params/eps_min}, thỉnh thoảng tác tử còn đi một bước
ngẫu nhiên. Khi tắt thăm dò (chính sách tham lam), tác tử đi đúng lộ trình tối ưu
\SL{ch8_qlearning/result/path_len} bước.

\section{Phân tích Q-Table và chính sách hành động tối ưu}
Bảng~\ref{tab:ch8-qtable} trích Q-Table tại ô xuất phát, các điểm rẽ của lộ trình và ô ngay trước đích; giá trị lớn nhất
mỗi hàng được in đậm và quyết định hành động $\pi^*(s)$.

\bang{ch8_qtable}{Trích đoạn Q-Table và chính sách hành động tối ưu $\pi^*(s)$ tại các ô then chốt.}{ch8-qtable}

Các giá trị trong bảng không phải con số tùy ý. Chúng hội tụ về đúng tổng phần thưởng chiết khấu mà phương trình Bellman
dự đoán, và có thể kiểm tra bằng tay:
\begin{itemize}
  \item \textbf{Ô ngay trước đích $(11, 10)$:} $Q(\text{Phải}) = \SL{ch8_qlearning/qtable/12/Q(Phải)}$, đúng bằng
        $R_{\text{goal}}$ vì bước này kết thúc episode, không còn phần thưởng tương lai.
  \item \textbf{Ô $(10, 10)$:} cả Phải và Xuống đều bằng \SL{ch8_qlearning/qtable/11/Q(Phải)} $= -1 + 0{,}95 \times
        100$. Hai hướng cùng dẫn tới một ô kề đích nên ngang giá trị; cả hai đều tối ưu.
  \item \textbf{Ô xuất phát $(0, 0)$:} $Q(\text{Phải}) = \SL{ch8_qlearning/qtable/0/Q(Phải)}$ là giá trị của lộ trình
        \SL{ch8_qlearning/result/path_len} bước:
        $\displaystyle -\sum_{k=0}^{22} 0{,}95^k + 0{,}95^{23} \times 100 \approx -13{,}85 + 30{,}73 = 16{,}88$.
  \item \textbf{Hành động đâm tường:} tại $(0,0)$, Trái và Lên đưa robot ra khỏi mê cung nên robot đứng yên:
        $Q = -10 + 0{,}95 \times 16{,}88 = \SL{ch8_qlearning/qtable/0/Q(Trái)}$. Giá trị này thấp hơn hẳn hướng tốt, nên
        tác tử đã học được cách tránh tường lửa.
\end{itemize}
Giá trị $Q$ dọc lộ trình tăng dần từ \SL{ch8_qlearning/qtable/0/Q(Phải)} ở $S$ lên 100 ở sát đích, mỗi bước nhân
thêm $\gamma$ và cộng phần thưởng bước. Q-Table vì thế đóng vai một ``bản đồ độ dốc'': từ bất kỳ ô nào trên lộ trình,
đi theo hướng có $Q$ lớn nhất là đi về đích nhanh nhất.
